\section{Efficiency of a memoryless transfer against a direct in-band drive}\label{app:bound}
For the logistic family of memoryless transfers, compare balanced TI tones of
polarization amplitude $\varepsilon$ each with a direct in-band drive of amplitude
$2\varepsilon$, matching peak polarization. Their leading drive amplitudes are
$|\Sig''|\varepsilon^2/2$ and $2\Sig'\varepsilon$, so
$\mathcal R_{\rm det}\simeq|\Sig''|\varepsilon/(4\Sig')$.
The weak-input formula concerns the logistic family. The step-mixture argument
below extends the exact carrier-averaged bound to bounded nondecreasing
memoryless transfers, with the same operating point, matched peak polarization,
and open-loop readout. It does not cover the dynamic detector of \S\ref{sec:dynamic}.

Writing $\zeta=\Sig/(2e_0)\in(0,1)$ gives
$\Sig'=2e_0\rho\zeta(1-\zeta)$ and
$\Sig''=2e_0\rho^2\zeta(1-\zeta)(1-2\zeta)$. Since
$|\mathrm{PEIX}|=1-4\zeta(1-\zeta)=(1-2\zeta)^2$, the leading ratio satisfies
\begin{equation}
\mathcal R_{\rm det}\simeq\frac{\rho\,\varepsilon}{4}\sqrt{|\mathrm{PEIX}|}
\;\le\;\frac{\rho\,\varepsilon}{4}.
\label{eq:etabound}
\end{equation}
It vanishes at the inflection, where the quadratic coefficient is zero, and
approaches its bound in the saturated tails. This maximizes a ratio rather than
an absolute response: both drives vanish in those tails. For the JR slope
$\rho=0.56$~mV$^{-1}$, the perturbative bound is $1.4\%$ at
$\varepsilon=0.1$~mV and $0.28\%$ at $\varepsilon=0.02$~mV.

Optimizing over the sigmoid slope gives a different limit from the weak-input bound.
The logistic transfer depends on voltage only through $\rho(v-v_0)$, so writing
\begin{equation}
u=\rho\,\varepsilon,\qquad \xi=\rho\,(v^*-v_0),
\label{eq:reduced}
\end{equation}
we have $\Sig(v^*+\varepsilon\psi(t))=2e_0\,L(\xi+u\,\psi(t))$ for any normalized waveform $\psi$,
with $L(\xi)=(1+e^{-\xi})^{-1}$. Every Fourier amplitude, and hence their ratio, is a function
of $(u,\xi)$ alone; the factor $2e_0$ cancels. Maximizing over $v^*$ and $\rho$ at fixed
$\varepsilon$ is therefore the same as maximizing over $(u,\xi)$, and the family supremum is
independent of field strength. Evaluating the exact cycle response rather than its
quadratic truncation and optimizing $\xi$ at each $u$, the ratio rises monotonically over the
range explored and converges to the hard-threshold limit
\begin{equation}
\sup_{\rho,\,v^*}\mathcal R_{\rm det}\approx0.3479 ,
\label{eq:etasup}
\end{equation}
approached as $\rho\to\infty$ at a threshold offset $\delta$ (the distance from the
operating point to the inflection) with $\delta/2\varepsilon=0.523$. There is no
finite optimal transfer width: narrower is always better, and reaching $95\%$ of the ceiling
requires $u\gtrsim6$, that is $1/\rho\lesssim\varepsilon/6$.

The limiting detector switches across its full range. At the optimum,
and at the illustrative fast-access polarization $\varepsilon\!\simeq\!0.1$~mV
of \S\ref{sec:discussion}, its open-loop beat drive
exceeds the physiological-slope sigmoid's by some three orders of magnitude for the same
$e_0$ and waveform, because the field excursion is about twice the threshold offset and the
transfer switches across its full range instead of responding linearly. The disparity grows
as the field weakens, the smooth sigmoid's response falling as $\varepsilon^2$ while a
threshold repositioned with the field keeps its excursion.

The required transfer width is small: at $\varepsilon=0.1$~mV the ceiling is approached only for
$1/\rho\lesssim17\;\mu$V; at $\varepsilon=0.02$~mV the corresponding width is
about $3\;\mu$V. Both follow from $1/\rho\lesssim\varepsilon/6$, with the
threshold offset held to the same order as the applied polarization. A population
sigmoid averages the spread of thresholds and noise; the JR voltage scale is
$1/\rho\approx1.8$~mV \cite{jansen1995}. A
microvolt width lies far below that and reads as an idealized microscopic threshold or
rectifier rather than as a neural-mass sigmoid; real thresholds are themselves noisy and
state-dependent, so it is not a measured single-cell property either. The
regime in which the family approaches its ceiling is thus one in which the ``population
sigmoid'' has ceased to describe a population.

The limit is a step detector, not an envelope follower. For a unit step at offset
$\delta$ the carrier-averaged transduction is $\pi^{-1}\arccos(\delta/\mathcal A)$ for $\mathcal A>\delta$
and zero below; with $\beta=\delta/2\varepsilon$ its tACS fundamental is
$(2/\pi)\sqrt{1-\beta^2}$ and its TI coefficient reduces by parts to
$\mathcal R_{\rm det}=2\beta\,[\mathsf K(k_{\rm ell}^2)-\mathsf E(k_{\rm ell}^2)]/(\pi\sqrt{1-\beta^2})$
with elliptic parameter $k_{\rm ell}^2=1-\beta^2$ and $\mathsf K,\mathsf E$
the complete elliptic integrals, giving $0.3479$ at $\beta=0.523$ by direct quadrature and
by the elliptic form alike (\texttt{threshold\_bound\_check.py}). At that optimum the
step's second-to-first harmonic magnitude ratio is $0.151$, not the $1/5$ of a linear
envelope follower. A bounded nondecreasing transfer is, up to a constant, a positive
mixture of shifted steps. Both signed Fourier coefficients are linear in the
transfer, and each step's direct-drive coefficient is nonnegative. The triangle
inequality therefore bounds the mixture's amplitude ratio by the largest step
ratio, whenever the denominator is nonzero. This extends the result beyond the
logistic family; the maximum itself is a numerical evaluation, not a certified bound.

\subsection{Experimental amplitude and phase-locking comparisons}\label{supp:vieira}
The protocols in Vieira et al.\ \cite{vieira2024} do not match local polarization at the detector.
The two-pair TI comparison gave a phase-locking effect ratio of $0.34$ against
tACS (95\% confidence interval $0.04$--$0.65$, 65 of 154 neurons), without
correcting for TI currents up to $\pm2.5$~mA per pair versus $\pm1$~mA for tACS.
A separate single-pair AM-tACS control compensated for frequency-dependent
shunting to match extracellular modulation depth. Of twenty cells, ten responded
to either modality and seven to both; in the latter group, median phase-locking
value (PLV) was $0.10$ for AM-tACS and $0.20$ for tACS, against $0.043$ at baseline.
The reported TI modulation depths were about $0.5$--$0.62$~V/m.

